% GENERATED FILE. Edit canonical lessons, then run npm run generate:books.
% canonical-source-hash: fnv1a64:7f0c3712b02c4300
% canonical-lessons: IT-C523-mostra, IT-C523-capolavoro, IT-C523-ritratto, IT-C523-affresco, IT-C523-restaurare

\chapter{Una mostra}
\label{ch:it-una-mostra}
\hlchaptermodality{523}

\begin{chapteropening}
\textbf{By the end of this chapter:} \emph{I can talk about visiting an exhibition.}

Complete the last lesson of chapter 523: I can talk about visiting an exhibition.
\end{chapteropening}

\section[la mostra]{la mostra — exhibition, show}

\label{lesson:IT-C523-mostra}



\begin{quote}
Before the new one: say the Italian for clay, then the Italian for to model.
\end{quote}

\subsection*{You'll want to know: la mostra}
\textbf{la mostra} — \textquotedblleft{}exhibition, show\textquotedblright{}.

\begin{quote}
Domenica andiamo a vedere una mostra di fotografia.
\end{quote}

\begin{quote}
\emph{On Sunday we are going to see a photography exhibition.}
\end{quote}

The same form is a verb: \emph{mi mostra la casa} — he shows me the house.

\begin{cousinweb}
\emph{mostra} is built on \emph{mostrare} \textquotedblleft{}to show\textquotedblright{}, inherited from Latin \emph{mōnstrāre} \textquotedblleft{}to point out\textquotedblright{} (EtymDB).
\end{cousinweb}

\begin{grammarlens}[title={What you've built}]
One more word for telling, explaining and describing at B1.
\end{grammarlens}

\subsection*{Your turn}
\begin{itemize}
\raggedright
  \item \emph{Say it:} \emph{la mostra}
  \item \emph{Say it:} \emph{la mostra}, once more
  \item \emph{Say it:} a sentence of your own with \emph{la mostra}
  \item \emph{Recall:} say the Italian for clay, then the Italian for to model, then say \emph{la mostra} again
\end{itemize}

\begin{tcolorbox}[breakable,colback=teal!4,colframe=teal!35!black,title={Before you move on}]
What does la mostra mean? (\textquotedblleft{}Exhibition, show\textquotedblright{}.) Say it once more.
\end{tcolorbox}
\section[il capolavoro]{il capolavoro — masterpiece}

\label{lesson:IT-C523-capolavoro}



\begin{quote}
Before the new one: say the Italian for to model, then the Italian for exhibition.
\end{quote}

\subsection*{You'll want to know: il capolavoro}
\textbf{il capolavoro} — \textquotedblleft{}masterpiece\textquotedblright{}.

\begin{quote}
Questo quadro è un capolavoro.
\end{quote}

\begin{quote}
\emph{This painting is a masterpiece.}
\end{quote}

Literally \textquotedblleft{}head work\textquotedblright{}: the best work.

\begin{cousinweb}
\emph{capolavoro} is built from \emph{capo} \textquotedblleft{}head\textquotedblright{}, from Latin \emph{caput} (Lewis \& Short), and \emph{lavoro} \textquotedblleft{}work\textquotedblright{}, from Latin \emph{labor}.
\end{cousinweb}

\begin{grammarlens}[title={What you've built}]
One more word for telling, explaining and describing at B1.
\end{grammarlens}

\subsection*{Your turn}
\begin{itemize}
\raggedright
  \item \emph{Say it:} \emph{il capolavoro}
  \item \emph{Say it:} \emph{il capolavoro}, once more
  \item \emph{Say it:} a sentence of your own with \emph{il capolavoro}
  \item \emph{Recall:} say the Italian for to model, then the Italian for exhibition, then say \emph{il capolavoro} again
\end{itemize}

\begin{tcolorbox}[breakable,colback=teal!4,colframe=teal!35!black,title={Before you move on}]
What does il capolavoro mean? (\textquotedblleft{}Masterpiece\textquotedblright{}.) Say it once more.
\end{tcolorbox}
\section[il ritratto]{il ritratto — portrait}

\label{lesson:IT-C523-ritratto}



\begin{quote}
Before the new one: say the Italian for exhibition, then the Italian for masterpiece.
\end{quote}

\subsection*{You'll want to know: il ritratto}
\textbf{il ritratto} — \textquotedblleft{}portrait\textquotedblright{}.

\begin{quote}
Nel salotto c'è un ritratto della nonna da giovane.
\end{quote}

\begin{quote}
\emph{In the living room there is a portrait of grandmother as a young woman.}
\end{quote}

Also a description: \emph{un ritratto della città}.

\begin{cousinweb}
\emph{ritratto} is the past participle of \emph{ritrarre} \textquotedblleft{}to portray\textquotedblright{}, from Latin \emph{retrahere} \textquotedblleft{}to draw back\textquotedblright{} (EtymDB), built on \emph{trahere} \textquotedblleft{}to draw\textquotedblright{}.
\end{cousinweb}

\begin{grammarlens}[title={What you've built}]
One more word for telling, explaining and describing at B1.
\end{grammarlens}

\subsection*{Your turn}
\begin{itemize}
\raggedright
  \item \emph{Say it:} \emph{il ritratto}
  \item \emph{Say it:} \emph{il ritratto}, once more
  \item \emph{Say it:} a sentence of your own with \emph{il ritratto}
  \item \emph{Recall:} say the Italian for exhibition, then the Italian for masterpiece, then say \emph{il ritratto} again
\end{itemize}

\begin{tcolorbox}[breakable,colback=teal!4,colframe=teal!35!black,title={Before you move on}]
What does il ritratto mean? (\textquotedblleft{}Portrait\textquotedblright{}.) Say it once more.
\end{tcolorbox}
\section[l'affresco]{l'affresco — fresco}

\label{lesson:IT-C523-affresco}



\begin{quote}
Before the new one: say the Italian for masterpiece, then the Italian for portrait.
\end{quote}

\subsection*{You'll want to know: l'affresco}
\textbf{l'affresco} — \textquotedblleft{}fresco\textquotedblright{}.

\begin{quote}
Nella chiesa ci sono affreschi molto antichi.
\end{quote}

\begin{quote}
\emph{In the church there are very old frescoes.}
\end{quote}

\emph{A fresco}: painted on wet plaster.

\begin{cousinweb}
\emph{affresco} is built from \emph{a fresco} \textquotedblleft{}on the fresh (plaster)\textquotedblright{}; \emph{fresco} comes from Medieval Latin \emph{friscus}, from a Germanic word \emph{friskaz} \textquotedblleft{}fresh\textquotedblright{} (EtymDB).
\end{cousinweb}

\begin{grammarlens}[title={What you've built}]
One more word for telling, explaining and describing at B1.
\end{grammarlens}

\subsection*{Your turn}
\begin{itemize}
\raggedright
  \item \emph{Say it:} \emph{l'affresco}
  \item \emph{Say it:} \emph{l'affresco}, once more
  \item \emph{Say it:} a sentence of your own with \emph{l'affresco}
  \item \emph{Recall:} say the Italian for masterpiece, then the Italian for portrait, then say \emph{l'affresco} again
\end{itemize}

\begin{tcolorbox}[breakable,colback=teal!4,colframe=teal!35!black,title={Before you move on}]
What does l'affresco mean? (\textquotedblleft{}Fresco\textquotedblright{}.) Say it once more.
\end{tcolorbox}
\section[restaurare]{restaurare — to restore}

\label{lesson:IT-C523-restaurare}



\begin{quote}
Before the new one: say the Italian for portrait, then the Italian for fresco.
\end{quote}

\subsection*{You'll want to know: restaurare}
\textbf{restaurare} — \textquotedblleft{}to restore\textquotedblright{}.

\begin{quote}
Stanno restaurando la facciata del museo.
\end{quote}

\begin{quote}
\emph{They are restoring the front of the museum.}
\end{quote}

The noun is \emph{il restauro}.

\begin{cousinweb}
\emph{restaurare} is a learned borrowing from Latin \emph{restaurāre} \textquotedblleft{}to renew, to rebuild\textquotedblright{} (EtymDB).
\end{cousinweb}

\begin{grammarlens}[title={What you've built}]
One more word for telling, explaining and describing at B1.
\end{grammarlens}

\subsection*{Your turn}
\begin{itemize}
\raggedright
  \item \emph{Say it:} \emph{restaurare}
  \item \emph{Say it:} \emph{restaurare}, once more
  \item \emph{Say it:} a sentence of your own with \emph{restaurare}
  \item \emph{Recall:} say the Italian for portrait, then the Italian for fresco, then say \emph{restaurare} again
\end{itemize}

\begin{tcolorbox}[breakable,colback=teal!4,colframe=teal!35!black,title={Before you move on}]
What does restaurare mean? (\textquotedblleft{}To restore\textquotedblright{}.) Say it once more.
\end{tcolorbox}
